\documentclass[prc]{revtex4} \usepackage[dvips]{graphicx}
\usepackage{mathrsfs} \usepackage{amsfonts} \usepackage{lscape}

\usepackage{epic,eepic} \usepackage{amsmath} \usepackage{amssymb}
\usepackage[dvips]{epsfig} \usepackage[T1]{fontenc}
\usepackage{hyperref} \usepackage{bezier} \usepackage{pstricks}
\usepackage{dcolumn}% Align table columns on decimal point
\usepackage{bm}% bold math
%\usepackage{braket}
\usepackage[dvips]{graphicx} \usepackage{pst-plot}

\newcommand{\One}{\hat{\mathbf{1}}} \newcommand{\eff}{\text{eff}}
\newcommand{\Heff}{\hat{H}_\text{eff}}
\newcommand{\Veff}{\hat{V}_\text{eff}}
\newcommand{\braket}[1]{\langle#1\rangle}
\newcommand{\Span}{\operatorname{sp}}
\newcommand{\tr}{\operatorname{trace}}
\newcommand{\diag}{\operatorname{diag}}
\newcommand{\bra}[1]{\left\langle #1 \right|}
\newcommand{\ket}[1]{\left| #1 \right\rangle} \newcommand{\element}[3]
           {\bra{#1}#2\ket{#3}}

\newcommand{\normord}[1]{ \left\{#1\right\} }

\usepackage{amsmath}
% Sub-items of an exercise are labelled (i), (ii), ...
\renewcommand{\labelenumii}{(\roman{enumii})}
%\newcommand{\be}{\begin{equation}}
%\newcommand{\ee}{\end{equation}}
%\newcommand{\OP}[1]{{\bf\widehat{#1}}}

\begin{document}

\begin{center}
{\LARGE\bf
First midterm FYS4480 – Quantum mechanics for many-particle systems, deadline October 26, at midnight. This midterm counts 25$\%$ of the final grade. Total score:100 points.}
\end{center}
\begin{center}
% List of all institutions:
\centerline{{\small Department of Physics, University of Oslo, Norway}}
\end{center}


We present a simplified Hamiltonian consisting of an unperturbed
Hamiltonian and a so-called pairing interaction term. It is a model
which to a large extent mimics some central features of atomic
nuclei, certain atoms and systems which exhibit superfluidity or
superconductivity.  To study this system we will use three of the
methods of this course, in the order in which they have been
developed in the lectures: full configuration interaction (FCI)
theory, Hartree-Fock (HF) theory and the density matrix
renormalization group (DMRG). The midterm is divided into three parts
which follow this order and build on each other.

In Part I we use second quantization to set up the Hamiltonian matrix
in the space of all Slater determinants of the model. Diagonalizing
this matrix is the FCI method and provides us with the exact answer,
which serves as the reference for everything that follows. We also
look at what happens when the space of Slater determinants is
truncated. In Part II we go to the opposite extreme and look for the
best possible \emph{single} Slater determinant, which is Hartree-Fock
theory. In Part III we return to the exact FCI state and rewrite it as
a matrix product state (MPS), and the Hamiltonian as a matrix product
operator (MPO). The bond dimension $\chi$ of the MPS interpolates
between a product state over the single-particle levels ($\chi=1$) and
the exact FCI state, and the DMRG is the variational method which
finds the best MPS at a given $\chi$. The three methods are thus
different ways of parametrizing the same many-body wave function, and
the aim of the midterm is to see how they are related for one and the
same model.  When setting up the Hamiltonian matrix you will need to
solve an eigenvalue problem.

We define first the Hamiltonian, with a definition of the model space
and the single-particle basis. Thereafter, we present the various
exercises.


The Hamiltonian acting in the complete Hilbert space (usually infinite
dimensional) consists of an unperturbed one-body part, $\hat{H}_0$,
and a perturbation $\hat{V}$.

We limit ourselves to at most two-body interactions, our Hamiltonian
is then represented by the following operators
\[
\hat{H} = \sum_{\alpha\beta}\langle \alpha |h_0|\beta\rangle
a_{\alpha}^{\dagger}a_{\beta}
+\frac{1}{4}\sum_{\alpha\beta\gamma\delta}\langle \alpha\beta|
V|\gamma\delta\rangle
a_{\alpha}^{\dagger}a_{\beta}^{\dagger}a_{\delta}a_{\gamma},
\]
where $a_{\alpha}^{\dagger}$ and $a_{\alpha}$ etc.~are standard
fermion creation and annihilation operators, respectively, and
$\alpha\beta\gamma\delta$ represent all possible single-particle
quantum numbers.  The full single-particle space is defined by the
completeness relation $\hat{{\bf 1}} = \sum_{\alpha
  =1}^{\infty}|\alpha \rangle \langle \alpha|$.  In our calculations
we will let the single-particle states $|\alpha\rangle$ be
eigenfunctions of the one-particle operator $\hat{h}_0$. Note that the two-body part of the Hamiltonian 
contains anti-symmetrized matrix elements.


The above Hamiltonian acts in turn on various many-body Slater
determinants constructed from the single-basis defined by the one-body
operator $\hat{h}_0$.  As an example, the two-particle model space
$\mathcal{P}$ is defined by an operator
\[
\hat{P} = \sum_{\alpha\beta =1}^{m}|\alpha\beta \rangle \langle
\alpha\beta|,
\]
where we assume that $m=\dim(\mathcal{P})$ and the full space is
defined by
\[
\hat{P}+\hat{Q}=\hat{{\bf 1}},
\]
with the projection operator
\[
\hat{Q} = \sum_{\alpha\beta =m+1}^{\infty}|\alpha\beta \rangle \langle
\alpha\beta|,
\]
being the complement of $\hat{P}$.


Our specific model consists of $N$ doubly-degenerate and equally
spaced single-particle levels labelled by $p=1,2,\dots$ and spin
$\sigma=\pm 1$.  These states are schematically portrayed in
Fig.~\ref{fig:schematic}.  The first two single-particle levels define
a possible model space, indicated by the label $\mathcal{P}$.  The
remaining states span the excluded space $\mathcal{Q}$.

We write the Hamiltonian as
\[ \hat{H} = \hat{H}_0 + \hat{V} , \]
where
\[
\hat{H}_0=\xi\sum_{p\sigma}(p-1)a_{p\sigma}^{\dagger}a_{p\sigma}
\]
and
\[
\hat{V}=-\frac{1}{2}g\sum_{pq}a^{\dagger}_{p+}
a^{\dagger}_{p-}a_{q-}a_{q+}.
\]
Here, $H_0$ is the unperturbed Hamiltonian with a spacing between
successive single-particle states given by $\xi$, which we will set to
a constant value $\xi=1$ without loss of generality (except in
exercise~7 below, where we study the degenerate case $\xi=0$). The two-body
operator $\hat{V}$ has one term only. It represents the pairing
contribution and carries a constant strength $g$. The pairing
interaction is attractive for $g>0$.

The indices
$\sigma=\pm$ represent the two possible spin values. The interaction
can only couple pairs and excites therefore only two particles at the
time, as indicated by the rightmost four-particle state in
Fig.~\ref{fig:schematic}. There one of the pairs is excited to the
state with $p=9$ and the other to the state $p=7$. The two middle
possibilities are not possible with the present model.  We label
single-particle states within the model space as hole-states. The
single-particle states outside the model space are then particle
states. 

In our model we have kept both the interaction strength and the
single-particle level as constants.  In a realistic system like an
atom or the atomic nucleus this is not the case.





\begin{figure*}[htbp]
\vspace{1.0cm} \setlength{\unitlength}{1cm}
 \begin{picture}(15,14)
 \thicklines \put(-0.6,1){\makebox(0,0){$p=1$}}
 \put(-0.6,2){\makebox(0,0){$p=2$}} \put(-0.6,3){\makebox(0,0){$p=3$}}
 \put(-0.6,4){\makebox(0,0){$p=4$}} \put(-0.6,5){\makebox(0,0){$p=5$}}
 \put(-0.6,6){\makebox(0,0){$p=6$}} \put(-0.6,7){\makebox(0,0){$p=7$}}
 \put(-0.6,8){\makebox(0,0){$p=8$}} \put(-0.6,9){\makebox(0,0){$p=9$}}
 \put(-0.6,10){\makebox(0,0){$p=10$}}
 \put(-0.6,11){\makebox(0,0){$p=\dots$}}
 \put(16,8){\makebox(0,0){$\mathcal{Q}$}}
 \put(16,2){\makebox(0,0){$\mathcal{P}$}}
% first 4-particle state
\put(0.8,1){\circle*{0.3}} \put(0.8,2){\circle*{0.3}}
\put(1.7,1){\circle*{0.3}} \put(1.7,2){\circle*{0.3}}
% second 4-particle state
\put(5.0,1){\circle*{0.3}} \put(5.9,1){\circle*{0.3}}
\put(5.0,4){\circle*{0.3}} \put(5.0,2){\circle{0.3}}
\put(5.9,3){\circle*{0.3}} \put(5.9,2){\circle{0.3}}
% third 4-particle state
\put(9.2,1){\circle*{0.3}} \put(10.1,3){\circle*{0.3}}
\put(9.2,4){\circle*{0.3}} \put(10.1,8){\circle*{0.3}}
\put(10.1,1){\circle{0.3}} \put(9.2,2){\circle{0.3}}
\put(10.1,2){\circle{0.3}}
% third 4-particle state
\put(13.4,7){\circle*{0.3}} \put(14.3,7){\circle*{0.3}}
\put(13.4,9){\circle*{0.3}} \put(14.3,9){\circle*{0.3}}
\put(13.4,1){\circle{0.3}} \put(14.3,1){\circle{0.3}}
\put(13.4,2){\circle{0.3}} \put(14.3,2){\circle{0.3}}
\dashline[+1]{2.5}(0,1)(15,1) \dashline[+1]{2.5}(0,2)(15,2)
\dashline[+1]{2.5}(0,3)(15,3) \dashline[+1]{2.5}(0,4)(15,4)
\dashline[+1]{2.5}(0,5)(15,5) \dashline[+1]{2.5}(0,6)(15,6)
\dashline[+1]{2.5}(0,7)(15,7) \dashline[+1]{2.5}(0,8)(15,8)
\dashline[+1]{2.5}(0,9)(15,9) \dashline[+1]{2.5}(0,10)(15,10)
\thinlines \dashline{0.1}(0,2.5)(15,2.5) \dashline{0.1}(0,11)(15,11)
 \end{picture}
\caption{Schematic plot of the possible single-particle levels with
  double degeneracy.  The filled circles indicate occupied particle
  states while the empty circles represent vacant particle(hole)
  states.  The spacing between each level $p$ is constant in this
  picture.  The first two single-particle levels define our possible
  model space, indicated by the label $\mathcal{P}$.  The remaining
  states span the excluded space $\mathcal{Q}$.  The first state to
  the left represents a possible ground state representation for a
  four-fermion system. In the second state to the left, one pair is
  broken. This possibility is however not included in our
  interaction. \label{fig:schematic}}
\end{figure*}

%======================================================================
\subsection*{Part I: Second quantization and full configuration interaction theory (40 points)}

\begin{enumerate}
\item (15/100 points) Show that the unperturbed Hamiltonian $\hat{H}_0$ and $\hat{V}$
  commute with both the spin projection $\hat{S}_z$ and the total spin
  $\hat{S}^2$, given by
\[
  \hat{S}_z := \frac{1}{2}\sum_{p\sigma} \sigma
  a^\dag_{p\sigma}a_{p\sigma}
\]
and
\[
  \hat{S}^2 := \hat{S}_z^2 + \frac{1}{2}(\hat{S}_+\hat{S}_- +
  \hat{S}_-\hat{S}_+),
\]
where
\[
  \hat{S}_\pm := \sum_{p} a^\dag_{p\pm} a_{p\mp}.
\]

This is an important feature of our system that allows us to
block-diagonalize the full Hamiltonian. We will focus on total spin
$S=0$.  In this case, it is convenient to define the so-called pair
creation and pair annihilation operators
\[
\hat{P}^{+}_p = a^\dag_{p+}a^\dag_{p-},
\]
and
\[
\hat{P}^{-}_p = a_{p-}a_{p+},
\] 
respectively.

Show that you can rewrite the Hamiltonian (with $\xi=1$) as
\[
\hat{H}=\sum_{p\sigma}(p-1)a_{p\sigma}^{\dagger}a_{p\sigma}
-\frac{1}{2}g\sum_{pq}\hat{P}^{+}_p\hat{P}^{-}_q.
\]
Show also that pair creation operators  commute among themselves.

In this midterm we focus only on a 
system with no broken pairs. This means that the Hamiltonian can only
link two-particle states in so-called spin-reversed states.


\item (15/100 points) Construct thereafter the Hamiltonian matrix for a system with no
  broken pairs and total spin $S=0$ for the case of the four lowest
  single-particle levels indicated in the
  Fig.~\ref{fig:schematic}. Our system consists of four particles
  only.  Our single-particle space consists of only the four lowest
  levels $p=1,2,3,4$.  You need to set up all possible Slater
  determinants.  Find all eigenvalues by diagonalizing the Hamiltonian
  matrix.  Vary your results for values of $g\in [-1,1]$.  We refer to
  this as the exact calculation. Comment the behavior of the ground
  state as function of $g$. Keep the ground-state eigenvector (not
  only the eigenvalue) for each value of $g$; you will need it in Part~III.

\item  (10/100 points)

Instead of setting up all possible Slater determinants, construct only
an approximation to the ground state (where we assume that the four
particles are in the two lowest single-particle orbits only) which
includes at most two-particle-two-hole excitations. Diagonalize this
matrix and compare with the exact calculation and comment your
results. Can you set up which diagrams this approximation corresponds
to?

\end{enumerate}

%======================================================================
\subsection*{Part II: Hartree-Fock theory (15 points)}

\begin{enumerate}
\setcounter{enumi}{3}
\item  (15/100 points) We switch now to approximative methods, with Hartree-Fock
  theory first. Hereafter we will define
  our model space to consist of the single-particle levels $p=1,2$.
  The remaining levels $p=3,4$ define our excluded space.  This means
  that our ground state Slater determinant consists of four particles
  which can be placed in the doubly degenerate orbits $p=1$ and $p=2$.
  Our first step is to perform a Hartree-Fock calculation with the
  pairing Hamiltonian.  Write first the normal-ordered Hamiltonian
  with respect to the above reference state given by four spin $1/2$
  fermions in the single-particle levels $p=1,2$, and set up the
  standard Hartree-Fock equations for the above system (often called
  restricted Hartree-Fock due to the fact that we have an equal number
  of spin-orbitals).

 The standard way to perform a Hartree-Fock calculation is to
expand the single-particle functions in a known basis  and vary the coefficients,  that is, the new function single-particle wave function $|p\rangle$ is written as a linear expansion in terms of a fixed basis $\phi$
\begin{equation*} 
\psi_p  = \sum_{\lambda} C_{p\lambda}\phi_{\lambda},
\end{equation*}

This leads to a new Slater determinant which is related to the
previous via a unitary transformation.  
We consider a Slater determinant built up of single-particle orbitals
$\phi_{\lambda}$ where the indices $\lambda$ refer to specific
single-particle states. 

The unitary transformation

\begin{equation*}
\psi_p  = \sum_{\lambda} C_{p\lambda}\phi_{\lambda},
\end{equation*}
brings us into the new basis $\psi$.  The new basis is orthonormal and $C$ is a unitary matrix.


Minimizing with respect to $C^*_{p\alpha}$, remembering that
$C^*_{p\alpha}$ and $C_{p\alpha}$ (and that the indices contain all
single-particle quantum numbers including spin) are independent and
defining

\begin{equation*}
h_{\alpha\gamma}^{HF}=\langle \alpha | h | \gamma \rangle+
\sum_{p}\sum_{\beta\delta} C^*_{p\beta}C_{p\delta}\langle \alpha\beta|V|\gamma\delta\rangle_{AS},
\end{equation*}
show that you can write the Hartree-Fock  equations as

\begin{equation*}
\sum_{\gamma}h_{\alpha\gamma}^{HF}C_{p\gamma}=\epsilon_p^{\mathrm{HF}}C_{p\alpha}.
\label{eq:newhf}
\end{equation*}

Explain the meaning of the different terms and define the Hartree-Fock
operator in second quantization. Write down its diagrammatic
representation as well.  The greek letters refer to the wave functions
in the original basis 
while roman letters refer to the new basis.


We will now set up the Hartree-Fock equations by varying the
coefficients of the single-particle functions. The single-particle
basis functions are defined as
\[
\psi_p = \sum_{\lambda} C_{p\lambda}\psi_{\lambda}.
\]
where in our case $p=1,2,3,4$ and $\lambda=1,2,3,4$, that is the first
four lowest single-particle orbits of Fig.~\ref{fig:schematic}.  Set
up the Hartree-Fock equations for this system by varying the
coefficients $C_{p\lambda}$ and solve them for values of $g\in
[-1,1]$.  Comment your results and compare with the exact
solution. Compute the total ground state energy  using a Hartree-Fock basis and comment your results.
Does the Hartree-Fock single-particle basis differ from the original
basis $\phi_\lambda$ for this interaction? Keep the Hartree-Fock
ground-state energy and the Hartree-Fock Slater determinant for the
comparison in Part~III.

\end{enumerate}

%======================================================================
\subsection*{Part III: From Slater determinants to matrix product states and the DMRG (45 points)}

In Part~I the exact ground state was obtained as a vector in the space
of Slater determinants with no broken pairs, and in Part~II it was
approximated by a single Slater determinant. The density matrix
renormalization group offers a third parametrization of the same
state, in which the amount of correlation that is kept is controlled by
a single number, the bond dimension $\chi$. In this part we rewrite
the pairing model in this language and connect it to the two previous
parts. The relevant background is the chapter on the density matrix
renormalization group in the lecture notes; the programs
\texttt{dmrg.py} in \texttt{doc/BookManybody/BookPrograms/chapterdmrg/}
and the notebook \texttt{doc/LectureNotes/dmrg.ipynb} may be used
freely, but everything asked for below can also be done with
\texttt{numpy} alone.

\paragraph*{Notation.}
We regard each level $p=1,\dots,L$ as a \emph{site} of a chain, with
the levels ordered by their energy. With no broken pairs a level is
either empty or occupied by a pair, so each site has a local
dimension $d=2$ with basis states $|0\rangle$ and $|1\rangle$. The
pair creation and annihilation operators of exercise~1 act on one site
only, $\hat{P}^+_p|0\rangle_p=|1\rangle_p$ and
$\hat{P}^-_p|1\rangle_p=|0\rangle_p$, and the number operator of the
level is $\hat{n}_p=\hat{n}_{p+}+\hat{n}_{p-}$. The Slater
determinants of exercise~2 are then product states
$|n_1n_2\cdots n_L\rangle$ with $n_p\in\{0,1\}$; for four levels and
two pairs they are
\[
|1100\rangle,\quad |1010\rangle,\quad |1001\rangle,\quad
|0110\rangle,\quad |0101\rangle,\quad |0011\rangle,
\]
where for example $|1100\rangle=\hat{P}^+_1\hat{P}^+_2|0\rangle$ is the
reference state of exercise~4. If broken pairs are allowed, each site
has instead $d=4$ states, $|0\rangle$, $|+\rangle$, $|-\rangle$ and
$|+-\rangle$.

Any state of the chain can be written as a \emph{matrix product state}
(MPS)
\[
|\Psi\rangle=\sum_{n_1\cdots n_L}
A^{[1]n_1}A^{[2]n_2}\cdots A^{[L]n_L}\,|n_1n_2\cdots n_L\rangle,
\]
where $A^{[p]n_p}$ is a $\chi_{p-1}\times\chi_p$ matrix for every value
of $n_p$, with $\chi_0=\chi_L=1$ so that the product is a number. The
$\chi_\ell$ are the \emph{bond dimensions}; the bond $\ell$ is the cut
between site $\ell$ and site $\ell+1$. The smallest possible
$\chi_\ell$ is the rank of the \emph{coefficient matrix} of the cut,
\[
C_{(n_1\cdots n_\ell),(n_{\ell+1}\cdots n_L)}
=\langle n_1\cdots n_L|\Psi\rangle ,
\]
whose singular value decomposition $C=U\Sigma V^\dagger$ is the Schmidt
decomposition $|\Psi\rangle=\sum_{k=1}^{\chi_\ell}\sigma_k
|\ell_k\rangle|r_k\rangle$ across that cut. The \emph{entanglement
entropy} of the cut is $S_\ell=-\sum_k\sigma_k^2\ln\sigma_k^2$ for a
normalized state, and an MPS is in \emph{left-canonical form} when
$\sum_{n}(A^{[p]n})^\dagger A^{[p]n}=\mathbf{1}$ for every site $p$,
which is what one obtains by performing the singular value
decompositions from left to right and absorbing $\Sigma V^\dagger$
into the next site. A product state, such as a single Slater
determinant in the level basis, has $\chi_\ell=1$ for every cut and
zero entanglement entropy.

In the same way an operator can be written as a \emph{matrix product
operator} (MPO)
\[
\hat{H}=\sum_{\{n\},\{n'\}} W^{[1]n_1n_1'}W^{[2]n_2n_2'}\cdots
W^{[L]n_Ln_L'}\,|n_1\cdots n_L\rangle\langle n_1'\cdots n_L'| ,
\]
where it is convenient to think of $W^{[p]}$ as a $D\times D$ matrix
whose entries are operators acting on site $p$; the first site
contributes only the first row and the last site only the last column,
so that the product $W^{[1]}W^{[2]}\cdots W^{[L]}$ is a single operator
on the chain. The number $D$ is the bond dimension of the MPO. It is
helpful to read the product as an automaton walking along the chain,
whose state keeps track of which operators have been placed and are
still waiting for their partner.

The DMRG is the variational minimization of
$\langle\Psi|\hat{H}|\Psi\rangle/\langle\Psi|\Psi\rangle$ over all MPS
with bond dimensions at most $\chi$, carried out one or two sites at a
time. With $\chi=1$ it searches among product states over the levels;
with $\chi$ large enough it reproduces the exact FCI state.

\begin{enumerate}
\setcounter{enumi}{4}

\item (10/100 points) \textbf{The Hamiltonian as a matrix product operator.}
Starting from the pair form of the Hamiltonian derived in exercise~1,
\[
\hat{H}=\sum_{p=1}^{L}\varepsilon_p\hat{n}_p
-\frac{g}{2}\sum_{pq}\hat{P}^+_p\hat{P}^-_q,
\qquad \varepsilon_p=\xi(p-1),
\]
describe the effect of each term when it acts on a state
$|n_1n_2\cdots n_L\rangle$. Restrict to zero seniority, that is $d=2$,
and write down the $2\times 2$ matrices representing $\hat{P}^+_p$,
$\hat{P}^-_p$ and $\hat{n}_p$ in the basis $\{|0\rangle,|1\rangle\}$ of
one site. Explain why the MPO bond dimension is $D=4$, that is,
identify the four states of the automaton that builds the
Hamiltonian, and explain why the bond dimension does not grow with the
number of levels $L$ even though the interaction couples every level to
every other level. Write out the MPO matrix $W^{[p]}$ for a general
level $p$ and multiply out the full operator for the case of four
levels, $W^{[1]}W^{[2]}W^{[3]}W^{[4]}$, where $W^{[1]}$ is the first row
and $W^{[4]}$ the last column of the general matrix. Check that every
term of $\hat{H}$ appears exactly once, and that the resulting
operator, restricted to the states with two pairs, reproduces the
$6\times 6$ Hamiltonian matrix of exercise~2. Finally, what would the
MPO be if each site has $d=4$ states $|0\rangle$, $|+\rangle$,
$|-\rangle$ and $|+-\rangle$?

\item (5/100 points) \textbf{Fermions.}
The particles are fermions. How does this affect the MPO? Recall that
if the $2L$ spin-orbitals are ordered along the chain, a fermion
operator acting on orbital $m$ carries the parity string
$\prod_{m'<m}(-1)^{\hat{n}_{m'}}$ of all orbitals before it
(Jordan-Wigner transformation). Discuss both the $d=2$ and the $d=4$
representation, and explain why the pairing Hamiltonian is a
particularly simple case. Which additional terms in a Hamiltonian,
for instance a one-body term
$a^\dagger_{p\sigma}a_{q\sigma}$ with $p\neq q$ which moves a single
particle between levels, would require the parity strings to appear
explicitly in the MPO, and how would they appear?

\item (15/100 points) \textbf{Matrix product states by hand: the degenerate limit.}
Set $\xi=0$, so that all levels are degenerate and the Hamiltonian
consists of the pairing interaction alone.
\begin{enumerate}
\item For four levels and two pairs, explain why the ground state is
  the equal-amplitude superposition of the six ways of placing two
  pairs,
\[
|\psi\rangle=\frac{1}{\sqrt{6}}\Big(|1100\rangle+|1010\rangle+|1001\rangle
+|0110\rangle+|0101\rangle+|0011\rangle\Big),
\]
  where $|0\rangle$ refers to an empty level and $|1\rangle$ to a
  paired level. What is its energy, and what are the other eigenvalues
  of the $6\times 6$ matrix in this limit? Compare with the
  $g$-dependence found in exercise~2 for $\xi=1$.
\item Write down the coefficient matrix $C$ for each of the three cuts
  of the chain, and construct the MPS in left-canonical form.
\item What is the entanglement entropy of the middle cut, and how does
  it compare with the largest value that cut could carry? What is the
  entanglement entropy of the Hartree-Fock state of exercise~4?
\item The middle bond carries $\chi=3$ states. Say in words what
  distinguishes them, and use that to write down an MPS of $3\times 3$
  matrices (except at the edges) that generalizes to a model of $L$
  levels and two pairs. Close to the edges some of the three bond
  states are never used. Which, and why?
\item What are the bond dimensions $\chi_\ell$ for each bond
  $\ell=1,\dots,L-1$ for $N$ pairs in $L$ degenerate levels? How do they
  compare with the number of Slater determinants $\binom{L}{N}$ that
  the FCI calculation of exercise~2 has to handle, and with the maximal
  bond dimension $\min(2^\ell,2^{L-\ell})$ of an arbitrary state of the
  chain?
\end{enumerate}

\item (15/100 points) \textbf{From FCI via Hartree-Fock to the DMRG.}
Return to $\xi=1$ and to the exact ground state of exercise~2 for
$g\in[-1,1]$, written as a vector in the basis of the six states
$|n_1n_2n_3n_4\rangle$ above.
\begin{enumerate}
\item For each of the three cuts, compute numerically the singular
  value decomposition of the coefficient matrix, the Schmidt weights
  $\sigma_k^2$ and the entanglement entropy $S_\ell$. Plot the entropy
  of the middle cut as a function of $g$ and compare with the value
  found in exercise~7 for $\xi=0$. Comment on the limit $g\to 0$ and on
  the difference between $g>0$ and $g<0$. Which cut carries the most
  entanglement, and why?
\item The Hartree-Fock state of exercise~4 is a single Slater
  determinant. Check whether it is a product state in the basis
  $|n_1n_2n_3n_4\rangle$, and if so give its bond dimensions. Truncate
  the Schmidt decomposition of the exact state across the middle cut
  to $\chi=1$ and to $\chi=2$ Schmidt states, renormalize, and compute
  the energy expectation value $\langle\hat{H}\rangle$ and the
  discarded weight $\sum_{k>\chi}\sigma_k^2$ of the truncated states as
  functions of $g$. Compare with the exact energy of exercise~2, the
  truncated-CI energy of exercise~3 and the Hartree-Fock energy of
  exercise~4. Explain in which sense the $\chi=1$ state and the
  Hartree-Fock state are both ``best product states'', and why they
  need not be the same state nor give the same energy.
\item Build the MPO of exercise~5 numerically as an array of shape
  $(D,D,d,d)$ and contract $W^{[1]}W^{[2]}W^{[3]}W^{[4]}$ to a
  $16\times 16$ matrix. Verify that, restricted to the two-pair
  states, it reproduces the Hamiltonian matrix of exercise~2 for a few
  values of $g$. Then compute $\langle\Psi|\hat{H}|\Psi\rangle$ for the
  truncated MPS of the previous point by contracting the MPO with the
  MPS, and check that you obtain the same energies as above.
\end{enumerate}
If time permits, run a two-site DMRG (your own, or
\texttt{dmrg.py}) for the four-level system with $\chi=1,2,3$ and
compare with the three previous points, and thereafter for a larger
system, for example $L=8$ levels and four pairs or $L=12$ levels and
six pairs, where the seniority-zero FCI matrix is still small enough
to diagonalize. Study the entanglement entropy across every bond as a
function of $g$ and comment on where along the chain the entanglement
is concentrated. This last part is not required for full score.

\end{enumerate}



 \end{document}
